\subsection{DEFINITION OF AN ALGEBRA}

DEFINITION 1. Let A be a commutative ring. An algebra over A (or an A-algebra, or simply an algebra, when no confusion is to be feared) is a set E with a structure defined by giving the following:

\begin{enumerate}
\def\labelenumi{(\arabic{enumi})}
\item
  an A-module structure on E;
\item
  an A-bilinear mapping (II, § 3, no. 5) of \(\mathbf{E} \times \mathbf{E}\) into \(\mathbf{E}\) .
\end{enumerate}

The A-bilinear mapping of \(\mathbf{E} \times \mathbf{E}\) into \(\mathbf{E}\) occurring in this definition is called the multiplication of the algebra \(\mathbf{E}\) ; it is usually denoted by \((x, y) \mapsto x.y\) , or simply \((x, y) \mapsto xy\) .

Let \((\alpha_{i})_{i\in I}\) and \((\beta_{j})_{j\in J}\) be two families of elements of A, of finite support (I, § 2, no. 1). Then, for all families \((x_{i})_{i\in I}\) and \((y_{j})_{j\in J}\) of elements of E, the general distributivity formula (I, § 3, no. 4)

\[
\left(\sum_ {i \in I} \alpha_ {i} x _ {i}\right) \left(\sum_ {j \in J} \beta_ {j} y _ {j}\right) = \sum_ {(i, j) \in I \times J} \left(\alpha_ {i} \beta_ {j}\right) \left(x _ {i} y _ {j}\right)\tag{1}
\]

holds; in particular

\[
(\alpha x) y = x (\alpha y) = \alpha (x y) \quad \text { for } \alpha \in \mathrm{A}, x \in \mathrm{E} \text { and } y \in \mathrm{E}.\tag{2}
\]

The bilinear mapping \((x,y)\mapsto yx\) of E × E into E and the A-module structure on E define an A-algebra structure on E, called opposite to the given algebra structure. The set E with this new structure is called the opposite algebra to the algebra E; it is often denoted by \(E^{0}\) . The A-algebra E is called commutative if it is identical with its opposite, in other words if multiplication in E is commutative. An isomorphism of E onto \(E^{0}\) is also called an anti-automorphism of the algebra E.

When multiplication in the algebra E is associative, E is called an associative A-algebra. When multiplication in E admits an identity element (necessarily unique (I, §2, no. 1)), this element is called the unit element of E and E is called a unital algebra.

Examples. (1) Every commutative ring A can be considered as an (associative and commutative) A-algebra.

\begin{enumerate}
\def\labelenumi{(\arabic{enumi})}
\setcounter{enumi}{1}
\item
  Let E be a pseudo-ring (I, §8, no. 1). Multiplication on E and the unique Z-module structure on E define on E an associative Z-algebra structure.
\item
  Let \(\mathbf{F}\) be a set and \(\mathbf{A}\) a commutative ring. The set \(\mathbf{A}^{\mathbf{F}}\) of all mappings of \(\mathbf{F}\) into \(\mathbf{A}\) , with the product ring structure (I, § 8, no. 10) and the product
\end{enumerate}

A-module structure (II, §1, no. 5) is an associative and commutative A-algebra.

\begin{enumerate}
\def\labelenumi{(\arabic{enumi})}
\setcounter{enumi}{3}
\tightlist
\item
  Let E be an A-algebra; the internal laws \((x,y)\mapsto xy+yx\) and \((x,y)\mapsto xy-yx\) define (with the A-module structure on E) two A-algebra structures on E, which are not in general associative; the first law
\end{enumerate}

\[
(x, y) \mapsto x y + y x
\]

is always commutative.

DEFINITION 2. Given two algebras E, E' over a commutative ring A, a homomorphism of E into E' is a mapping \(f: \mathbf{E} \to \mathbf{E}'\) such that

\begin{enumerate}
\def\labelenumi{(\arabic{enumi})}
\item
  \(f\) is an A-module homomorphism;
\item
  \(f(xy) = f(x)f(y)\) for all \(x \in \mathbf{E}\) and \(y \in \mathbf{E}\) .
\end{enumerate}

Clearly the composition of two A-algebra homomorphisms is an A-algebra homomorphism. Every bijective algebra homomorphism is an isomorphism. Therefore A-algebra homomorphisms may be taken as morphisms of the species of A-algebra structure (Set Theory, IV, 2, no. 1). We shall always suppose in what follows that this choice of morphisms has been made. If E, E' are two A-algebras, let \(\operatorname{Hom}_{\mathbf{A}-\operatorname{alg}}(\mathbf{E}, \mathbf{E}')\) denote the set of A-algebra homomorphisms of E into E'.

Let E, E' be two algebras each with a unit element. A homomorphism of E into E' mapping the unit element of E to the unit element of E' is called a unital homomorphism (or unital algebra morphism).

